% !PRINTED 59
\weberpage{59}

%<<CONT
heit congruent ist, so ist $e$ nach (1) ein Theiler von $2^{\lambda-2}$, also eine Potenz von 2, und wenn $e < 2^{\lambda-2}$ wäre, so müsste

\[
5^{2^{\lambda-3}} \equiv 1 \pmod{2^\lambda}
\]

sein. Dies ist aber nicht möglich, denn es gilt für jedes $\lambda$, was grösser als 2 ist, die Formel

\[
5^{2^{\lambda-3}} = 1 + 2^{\lambda-1} h\ ^{1)}
\]

mit ungeradem $h$, eine Formel, die sich ebenso wie die Formel (1) durch vollständige Induction beweisen lässt. Es sind also, wenn $\lambda \geqq 3$ ist,
die $2^{\lambda-2}$ Potenzen

\weq{2}{1, 5, 5^2 \ldots 5^{2^{\lambda-2}-1}}

nach dem Modul $2^\lambda$ alle von einander verschieden. Nun ist eine Relation von der Form $5^h \equiv -5^\varkappa$ für den Modul 4, also um so mehr für jede höhere Potenz von 2, unmöglich, und folglich sind die Grössen

\weq{3}{-1, -5, -5^2 \ldots -5^{2^{\lambda-2}-1}}

nach dem Modul $2^\lambda$ sowohl unter einander als von den Grössen (2) verschieden, und da ihre gesammte Anzahl $2^{\lambda-1}$ beträgt, so ist jede ungerade Zahl einer und nur einer der Grössen (2), (3) nach dem Modul $2^\lambda$ congruent.

Setzen wir also

\weq{4}{a = 2, \quad b = 2^{\lambda-2},}

so erhalten wir folgenden Satz für $\lambda \geqq 3$:

\noindent\textbf{3.} \gesperrt{Für jede ungerade Zahl} $n$ \gesperrt{lässt sich ein nach dem Modulpaar} $a, b$ \gesperrt{völlig bestimmtes Zahlenpaar} $\alpha, \beta$ \gesperrt{angeben, so dass}

\[
n \equiv (-1)^\alpha 5^\beta \pmod{2^\lambda}
\]

\gesperrt{wird.}

Der Fall $\lambda = 2$ kann hierunter mit subsumirt werden, weil dann $b = 1$ wird und $\beta = 0$ gesetzt werden kann. Um auch den Fall $\lambda = 1$ mit darunter zu begreifen, der aber kein besonderes Interesse bietet, müsste man $a = 1, b = 1$ setzen.

\noindent\rule{4cm}{0.4pt}

\noindent $^{1)}$ Wollte man an Stelle der Zahl 5 die Zahl 3 als Basis nehmen, so würde diese Formel für $\lambda = 3$ noch nicht gültig sein, wohl aber für jedes grössere $\lambda$, und daher könnte 3 als Basis ebenso gut dienen, wie 5. Der Grund für die Bevorzugung der Basis 5 liegt darin, dass alle Potenzen dieser Basis $\equiv 1 \pmod{4}$ sind.
